\section{有限分辨率读出接口：效应、POVM 与粗粒化}

\subsection{效应语言：把“有限读出”抽象成正算子分解}

\begin{definition}[POVM 读出（接口定义）]
\label{def:povm}
令 $\mathcal H$ 为有限维复 Hilbert 空间。一个有限分辨率读出由一族效应 $\{E_k\}_{k\in K}$ 给出，满足
$$
E_k\succeq 0,\qquad \sum_{k\in K}E_k=I.
$$
对任意密度算子 $\rho\succeq 0,\ \Tr(\rho)=1$，定义读出概率
$$
P_k:=\Tr(\rho E_k).
$$
\end{definition}

\begin{remark}
Definition \ref{def:povm} 作为标准接口语言可被视为对“有限读出事件”的最一般封装。
其数学基础与稀释实现（Naimark/Stinespring、Kraus 表示）是量子信息中的标准定理；本文仅引用之，不将其作为折叠闭层推理的前提。
\end{remark}

\subsection{与纤维计数的对应}

在本文的主线中，粗结果分布 $p_m$ 来自纤维计数。由第 \ref{thm:born_as_counting_equivalence} 定理可知：在 $\mathcal H_m=\CC^{X_m}$ 上取
$$
\rho_m:=|\psi_m\rangle\langle\psi_m|,
\qquad
E_x:=|x\rangle\langle x|,
$$
则 $\Tr(\rho_m E_x)=p_m(x)$。

\begin{remark}
在更一般的读出语义中，一个“粗标签”可能对应多个基事件的合并（粗粒化）。这对应将投影族按分块求和得到的 POVM：
若 $K=\bigsqcup_{a\in A}K_a$ 为分块，则令
$
E_a:=\sum_{k\in K_a}E_k
$
仍满足 $\sum_{a\in A}E_a=I$。
这与折叠/残差结构中“进一步商化得到更粗结果”的形式一致。
\end{remark}

